% =====================================================================
%  Praesentation 1: DevContainer Template
%  Kompilieren:  pdflatex -> pdflatex   (oder latexmk -pdf)
%  Benoetigt hhntheme.sty im selben Ordner.
%  Autor/Institut unten ggf. anpassen.
% =====================================================================
\documentclass[aspectratio=169,11pt]{beamer}

\usepackage[utf8]{inputenc}
\usepackage[T1]{fontenc}
\usepackage[ngerman]{babel}
\usepackage{hhntheme}

% --- Metadaten -------------------------------------------------------
\renewcommand{\hhnkicker}{Forschungsinfrastruktur \textperiodcentered\ Template-Serie 1/2}
\title{Reproduzierbare Forschungsumgebungen}
\subtitle{Ein GPU-fähiges Dev-Container-Template für die Arbeitsgruppe}
\author{Christian Heinzmann}
\institute{Hochschule Heilbronn \textperiodcentered\ Promotionskolleg}
\date{\today}

% Kurzfassungen fuer die Fusszeile
\title[DevContainer-Template]{Reproduzierbare Forschungsumgebungen}
\author[C. Heinzmann]{Christian Heinzmann}

\begin{document}

% --- Titel -----------------------------------------------------------
\hhntitleframe

% --- Agenda ----------------------------------------------------------
\begin{frame}{Agenda}
  \tableofcontents
\end{frame}

% =====================================================================
\section{Motivation: Das Reproduzierbarkeitsproblem}
% =====================================================================
\begin{frame}{Warum die Umgebung Teil der Wissenschaft ist}
  \begin{itemize}
    \item \hl{,,Works on my machine``} --- Ergebnisse hängen an undokumentierten
          lokalen Zuständen (Treiber, Bibliotheksversionen, Pfade).
    \item \hl{Environment-Drift}: CUDA, Python und Paketversionen divergieren
          zwischen Personen und Maschinen --- oft unbemerkt.
    \item \hl{Onboarding-Kosten}: Tage, bis ein neues Mitglied überhaupt
          rechnen kann; jede Maschine ein Einzelfall.
    \item \hl{Heterogene Hardware}: Laptop, Workstation, geteilte GPU-Server ---
          dieselbe Pipeline, unterschiedliche Realität.
  \end{itemize}
  \vfill
  \begin{takeaway}
    Computational Reproducibility ist kein Nice-to-have, sondern Teil
    wissenschaftlicher Sorgfalt. Die \hl{Umgebung muss versioniert, geteilt und
    identisch instanziierbar} sein --- genau wie der Code selbst.
  \end{takeaway}
\end{frame}

% =====================================================================
\section{Überblick: Was das Template liefert}
% =====================================================================
\begin{frame}[fragile]{Das Versprechen: Klonen $\rightarrow$ Reopen $\rightarrow$ Coden}
  \begin{columns}[T]
    \begin{column}{0.52\textwidth}
      \textbf{Drei Schritte, plattformunabhängig:}
      \begin{enumerate}
        \item Repository klonen (inkl. Submodule)
        \item In VS Code \emph{,,Reopen in Container``}
        \item Loslegen --- Umgebung ist fertig
      \end{enumerate}
      \vspace{0.4em}
      \textbf{Host-Voraussetzungen (das ist alles):}
      \begin{itemize}
        \item Git \textperiodcentered\ Docker (Compose V2)
        \item VS Code + Dev-Containers-Extension
        \item NVIDIA-GPU + Treiber \emph{(optional)}
      \end{itemize}
    \end{column}
    \begin{column}{0.48\textwidth}
\begin{lstlisting}[language=myshell]
# 1) Klonen inkl. Submodul
git clone <repo-url>
git submodule update --init --recursive

# 2) VS Code: F1 ->
#    "Dev Containers: Reopen in Container"

# 3) fertig - identisch auf
#    Windows / macOS / Linux
\end{lstlisting}
      \begin{takeaway}[Eine Spezifikation]
        Eine deklarative Definition erzeugt auf \hl{jedem} Host dieselbe
        GPU-fähige, KI-unterstützte Arbeitsumgebung.
      \end{takeaway}
    \end{column}
  \end{columns}
\end{frame}

\begin{frame}{Was vorinstalliert ist}
  \begin{columns}[T]
    \begin{column}{0.5\textwidth}
      \textbf{Stack}
      \begin{itemize}
        \item Ubuntu 24.04 \textperiodcentered\ CUDA 12.6.2 + cuDNN
        \item Python 3.12 in isoliertem venv (\code{/app/venv})
        \item Node.js 20 (für die Claude-Code-CLI)
      \end{itemize}
      \vspace{0.3em}
      \textbf{KI-Werkzeuge ab Werk}
      \begin{itemize}
        \item Claude Code (CLI + VS-Code-Extension)
        \item GitHub Copilot + Copilot Chat
      \end{itemize}
    \end{column}
    \begin{column}{0.5\textwidth}
      \textbf{Editor-Tooling}
      \begin{itemize}
        \item $\sim$30 VS-Code-Extensions automatisch
        \item Python/Pylance/Debugpy, vollständige
              Jupyter-Suite, Docker, Draw.io, Git-Graph
        \item Formatierung \& Linting konfiguriert
              (100-Zeichen-Lineal als Coding-Rule)
      \end{itemize}
      \vspace{0.3em}
      \begin{takeaway}[Ziel]
        Der erste Commit, nicht das Setup, ist der erste produktive Schritt.
      \end{takeaway}
    \end{column}
  \end{columns}
\end{frame}

% =====================================================================
\section{Architektur: Drei Konfigurationsebenen}
% =====================================================================
\begin{frame}{Trennung der Zuständigkeiten}
  \begin{center}
  \begin{tikzpicture}[font=\small,every node/.style={align=center}]
    \tikzset{
      layer/.style={rounded corners=3pt,minimum width=10.6cm,minimum height=1.35cm,
                    text width=10.2cm,draw=hhnblue,line width=1pt},
    }
    \node[layer,fill=cardbg] (dc) at (0,0)
      {\textbf{\color{hhnblue}devcontainer.json} \;--\; IDE-Glue\\
       \scriptsize Extensions \textperiodcentered\ Default-Interpreter \textperiodcentered\ Feature \emph{claude-code} \textperiodcentered\ \code{postStartCommand}};
    \node[layer,fill=lightteal!25] (cp) at (0,-1.6)
      {\textbf{\color{hhnpetrol}docker-compose.yml} \;--\; Laufzeit\\
       \scriptsize GPU-Reservierung \textperiodcentered\ Mounts \textperiodcentered\ Ports \textperiodcentered\ \code{shm\_size}};
    \node[layer,fill=hhnbluelight!40] (df) at (0,-3.2)
      {\textbf{\color{hhnblue}Dockerfile} \;--\; Build-Zeit (unveränderliche Basis)\\
       \scriptsize OS \textperiodcentered\ System-Pakete \textperiodcentered\ CUDA/cuDNN \textperiodcentered\ venv};
    \draw[-{Stealth[length=2.5mm]},hhncyan,line width=1.2pt] (df.north) -- (cp.south);
    \draw[-{Stealth[length=2.5mm]},hhncyan,line width=1.2pt] (cp.north) -- (dc.south);
  \end{tikzpicture}
  \end{center}
  \vspace{-0.2em}
  {\footnotesize \hl{Designprinzip:} schwere System-Pakete werden ins Image
  \emph{gebacken} (Build-Zeit), die Python-Abhängigkeiten bleiben zur
  \emph{Laufzeit} die lebende Single Source of Truth.}
\end{frame}

% =====================================================================
\section{Reproduzierbarkeit \& GPU}
% =====================================================================
\begin{frame}[fragile]{Reproduzierbarkeit beginnt am OS-Layer}
  \begin{columns}[T]
    \begin{column}{0.5\textwidth}
      \begin{itemize}
        \item Basis-Image \hl{exakt gepinnt} --- kein \code{latest},
              sondern ein fixer Tag (rechts):
        \item Python systemseitig, aber in venv \code{/app/venv} isoliert ---
              überall als Default-Interpreter gesetzt (Dockerfile,
              devcontainer.json, settings.json, Workspace).
        \item Node.js 20 \emph{ausschließlich} für die Claude-Code-CLI.
      \end{itemize}
      \begin{takeaway}[Konsequenz]
        Jede Person erhält denselben CUDA\,/\,cuDNN\,/\,Python-Stack ---
        Reproduzierbarkeit startet am OS-/Treiber-Layer, nicht erst bei
        \code{pip}.
      \end{takeaway}
    \end{column}
    \begin{column}{0.5\textwidth}
\begin{lstlisting}[language=mydocker]
FROM nvidia/cuda:12.6.2-cudnn-devel-ubuntu24.04

RUN apt-get update && apt-get install -y \
    python3.12 python3-pip python3.12-venv \
    git curl docker.io docker-compose-v2 \
    poppler-utils chromium ffmpeg ...

# venv als Default-Interpreter
RUN python3 -m venv /app/venv
ENV VIRTUAL_ENV=/app/venv
ENV PATH="/app/venv/bin:$PATH"
\end{lstlisting}
    \end{column}
  \end{columns}
\end{frame}

\begin{frame}[fragile]{GPU-Defaults für ML-Workloads}
  \begin{columns}[T]
    \begin{column}{0.52\textwidth}
      \begin{itemize}
        \item \hl{Alle GPUs reserviert} über Compose
              \code{deploy.resources.reservations}.
        \item \hl{\code{shm\_size: 32\,GB}} --- löst den klassischen
              Shared-Memory-Engpass beim Daten-Loading.
        \item X11 durchgereicht (GUI-Apps), Port 4010 für Streamlit
              forwarded.
        \item Verifikation \emph{vor} jeder Projektarbeit:
      \end{itemize}
\begin{lstlisting}[language=myshell]
docker run --rm --gpus all \
  nvidia/cuda:12.6.2-base-ubuntu24.04 nvidia-smi
\end{lstlisting}
    \end{column}
    \begin{column}{0.48\textwidth}
\begin{lstlisting}[language=myyaml]
deploy:
  resources:
    reservations:
      devices:
        - driver: nvidia
          count: all
          capabilities: [gpu]
shm_size: '32gb'
volumes:
  - ${localWorkspaceFolder:-..}:/app/...
  - ~/.gitconfig:/etc/gitconfig:ro
  - ~/.ssh:/root/.ssh:ro
\end{lstlisting}
    \end{column}
  \end{columns}
\end{frame}

\begin{frame}[fragile]{Automation: idempotenter Start}
  \begin{itemize}
    \item \code{postStartCommand} ruft bei jedem Start \code{setup\_startup.sh}
          (mit \code{set -e}).
    \item Schritte: venv aktivieren, \code{requirements} installieren, X11
          freigeben, GitHub-SSH prüfen.
    \item \hl{Non-fatal by design}: jeder optionale Schritt ist mit
          \code{|| true} abgesichert --- ein fehlendes X-Display oder
          SSH-Key blockiert den Container nie.
    \item Credentials werden \hl{wiederverwendet statt eingebrannt}:
          \code{\textasciitilde/.gitconfig} und \code{\textasciitilde/.ssh}
          read-only gemountet.
  \end{itemize}
\begin{lstlisting}[language=myshell]
source /app/venv/bin/activate
python -m pip install --no-cache-dir -r .../requirements.txt
xhost +local:docker >/dev/null 2>&1 || true          # GUI, best effort
ssh -o StrictHostKeyChecking=accept-new -T git@github.com || true   # nur Info
\end{lstlisting}
\end{frame}

% =====================================================================
\section{KI-Werkzeuge als Infrastruktur}
% =====================================================================
\begin{frame}{KI-Assistenten sind first-class Umgebung}
  \begin{itemize}
    \item Claude Code wird \hl{dreifach} bereitgestellt:
          devcontainer-Feature \code{anthropics/.../claude-code:1},
          VS-Code-Extension und Node 20 im Image.
    \item \hl{Ein Regelsatz für beide Assistenten}: \code{CLAUDE.md},
          \code{.claude/settings.json} und \code{.github/copilot-instructions.md}
          verweisen auf \emph{dieselbe} \code{CodingRules.md}.
    \item Single Source of Truth: Konventionen einmal pflegen ---
          Claude Code \emph{und} Copilot folgen automatisch.
  \end{itemize}
  \vfill
  \begin{takeaway}[Brücke]
    Die KI-Werkzeuge sind kein Add-on, sondern Teil der reproduzierbaren
    Umgebung. \hl{Details in Präsentation 2} (Claude Project Template).
  \end{takeaway}
\end{frame}

% =====================================================================
\section{Verteiltes Arbeiten: Laptop \texorpdfstring{$\leftrightarrow$}{<->} GPU-Server}
% =====================================================================
\begin{frame}[fragile]{Denselben Container überall ausführen}
  \begin{columns}[T]
    \begin{column}{0.5\textwidth}
      \textbf{Variante A --- alles remote}\\
      Workspace \emph{und} Container auf dem GPU-Server (Remote-SSH +
      ,,Reopen in Container``).
      \vspace{0.5em}

      \textbf{Variante B --- Code lokal, GPU remote}\\
      Docker-Context über SSH: der Laptop behält den Code, der Container
      läuft auf dem Server --- \hl{GPU ausleihen, ohne Code zu kopieren}.
\begin{lstlisting}[language=myshell]
docker context create gpu \
  --docker host=ssh://gpu-server
docker context use gpu
\end{lstlisting}
    \end{column}
    \begin{column}{0.5\textwidth}
      \textbf{Begleitende Host-Playbooks (\code{docs/help/NewPC/})}
      \begin{itemize}
        \item \hl{Multi-Remote-Git}: ein \code{git push all main}
              $\rightarrow$ GitHub \emph{und} Uni-GitLab gleichzeitig.
        \item \hl{eduVPN Split-Tunneling}: nur das Campus-Subnetz durch
              den Tunnel (systemd-persistent).
        \item \hl{SMB/CIFS-Auto-Mount}: \code{vers=3.0}, \code{nofail},
              \code{x-systemd.automount}.
      \end{itemize}
      \begin{takeaway}
        Das Template adressiert die reale Topologie einer Arbeitsgruppe ---
        nicht nur die ideale Einzelmaschine.
      \end{takeaway}
    \end{column}
  \end{columns}
\end{frame}

% =====================================================================
\section{Kritische Einordnung \& Roadmap}
% =====================================================================
\begin{frame}[shrink]{Ehrliche Bilanz}
  \begin{columns}[T]
    \begin{column}{0.46\textwidth}
      \good{Stärken}
      \begin{itemize}
        \item Basis-Image am OS-Layer gepinnt
        \item GPU + ML-Defaults out-of-the-box
        \item KI-Tooling integriert \& geteilt
        \item Remote-/Multi-Maschinen-fähig
        \item Umfangreich dokumentiert
      \end{itemize}
    \end{column}
    \begin{column}{0.54\textwidth}
      \warn{Offene Punkte (Roadmap)}
      \begin{itemize}
        \item \warn{Python-Deps nicht gelockt} (nur
              \code{streamlit==1.42.2}) $\rightarrow$ Lockfile
              (\code{uv}/\code{pip-tools}, Hashes).
        \item \warn{Keine CI} $\rightarrow$ Actions: Image-Build,
              \code{black}/\code{pylint}, \code{pytest}.
        \item \warn{Doku-Drift}: ReadMe nennt CUDA 12.6.0/Ubuntu 22.04 \&
              \code{python:3.9}; Submodul \code{src/Template} vs.\ Doku
              \code{src/NewWandB}.
        \item \warn{Kein CPU-Fallback}: \code{count: all} hart kodiert;
              Doku verweist auf nicht existente \code{\#\,NVIDIA}-Marker.
        \item \warn{Sicherheit auf geteilten Servern}: \code{docker.sock} +
              ganzes \code{\textasciitilde/.ssh} gemountet, root im Container.
      \end{itemize}
    \end{column}
  \end{columns}
\end{frame}

\begin{frame}{Konkrete nächste Schritte}
  \begin{enumerate}
    \item \hl{Lockfile einführen} --- repräsentativen Dev-Stack
          (numpy, pandas, torch passend zu CUDA 12.6, pytest/black/pylint)
          vollständig pinnen, mit Hashes.
    \item \hl{Minimale CI} --- Image bauen, Format/Lint/Tests laufen lassen,
          \code{main} mit Required-Checks schützen (macht die ,,getestet vor
          Merge``-Regel real statt aspirativ).
    \item \hl{Konsistenz-Pass} --- ReadMes aus den echten Dateien speisen,
          Submodul-Namen vereinheitlichen, \code{.claude/worktrees/} aus der
          Versionierung nehmen.
    \item \hl{CPU-Profil} --- Compose-Profile oder
          \code{docker-compose.cpu.yml}, damit auch reine Laptops starten.
    \item \hl{Härtungsprofil} dokumentieren --- non-root-User mit Host-UID,
          SSH-Agent-Forwarding statt \code{\textasciitilde/.ssh}-Mount,
          \code{docker.sock} opt-in.
    \item \hl{FAIR-Anker} --- \code{LICENSE}, \code{CITATION.cff},
          Zenodo-DOI: vom Setup zum \emph{zitierbaren} Artefakt.
  \end{enumerate}
\end{frame}

% =====================================================================
\section{Zusammenfassung}
% =====================================================================
\begin{frame}{Zusammenfassung}
  \begin{itemize}
    \item Die \hl{Umgebung wird zum versionierten Artefakt} ---
          Reproduzierbarkeit ab dem OS-Layer.
    \item \hl{Ein Klick} zur GPU-fähigen, KI-unterstützten Arbeitsumgebung,
          identisch auf jedem Host.
    \item \hl{Skaliert} über Laptops und geteilte GPU-Server hinweg.
    \item Roadmap: \emph{Lockfile + CI + Doku-Konsolidierung + DOI} ---
          von ,,reproduzierbar in der Absicht`` zu ,,maschinell verifiziert``.
  \end{itemize}
  \vfill
  \begin{center}
    {\large\color{hhnblue}\textbf{Fragen \& Diskussion}}\\[2pt]
    {\footnotesize\color{darkgray} Fortsetzung: Präsentation 2 --- Claude Project Template}
  \end{center}
\end{frame}

\end{document}
